\section*{अध्याय \ref{ch:g9:landforms} --- स्थलरूप}

\begin{solution}{exo:g9:landforms:1}
(क) हवा। (ख) \omterm{def:g3:stones-pebbles-sand:rock}{चट्टान} को घोलता पानी (और \omterm{def:g9:landforms:karst}{गुफा} में कैल्साइट जमा करता पानी)। (ग) \omterm{def:g2:water-on-earth:ocean}{समुद्र}
(\omterm{def:g9:landforms:beach}{तटीय अपवाह})। (घ) बर्फ़। (ङ) \omterm{def:g2:water-on-earth:ocean}{समुद्र}। (च) \omterm{def:g1:landscapes-around-us:water}{नदी}।
\end{solution}

\begin{solution}{exo:g9:landforms:2}
(क) \omterm{def:g9:landforms:uplift}{उत्थान}। (ख) \omterm{def:g9:landforms:dune}{फिसलन फलक}। (ग) \omterm{def:g9:landforms:karst}{घोलरंध्र}। (घ) \omterm{def:g9:landforms:beach}{तटीय अपवाह}।
(ङ) \omterm{def:g9:landforms:karst}{कार्स्ट}।
\end{solution}

\begin{solution}{exo:g9:landforms:3}
(क) ग़लत: भुजाएँ हवा की दिशा में आगे की ओर इशारा करती हैं; गोल नाक हवा की ओर होती है।
(ख) सही। (ग) ग़लत: हवा उसे कण-दर-कण चलाती है, हल्की ढाल पर ऊपर और \omterm{def:g9:landforms:dune}{फिसलन फलक} से
नीचे। (घ) सही।
\end{solution}

\begin{solution}{exo:g9:landforms:4}
(ख), (घ), (ङ), (क), (ग)।
\end{solution}

\begin{solution}{exo:g9:landforms:5}
(क) पश्चिम से: भुजाएँ हवा की दिशा में आगे की ओर होती हैं। (ख) पूर्व की ओर, भुजाओं के बीच।
(ग) $60 \div 6 = 10$ वर्ष। (घ) बड़े \omterm{def:g9:landforms:dune}{टिब्बे} में ज़्यादा \omterm{def:g3:stones-pebbles-sand:sand}{रेत} होती है, और \omterm{def:g9:landforms:dune}{टिब्बे} के आगे
बढ़ने के लिए उस सारी \omterm{def:g3:stones-pebbles-sand:sand}{रेत} को शिखा के पार ले जाना पड़ता
है।
\end{solution}

\begin{solution}{exo:g9:landforms:6}
(क) $7 \times 1 = \qty{7}{m}$। (ख) $2 \div 12 = \frac{1}{6}$ वर्ष।
(ग) एक वर्ष के $\frac{1}{6}$ में सात वर्षों का पीछे हटना: $7 \times 6 = 42$
गुना तेज़। (घ) लहरें पहले एक खाँच काटती हैं; ऊपर की \omterm{def:g3:stones-pebbles-sand:rock}{चट्टान} कुछ समय तक उस पर लटकी रहती है,
बिना हिले; फिर लटका हुआ हिस्सा एक साथ गिर जाता है, और किनारा छलाँग लगाकर पीछे हट जाता
है।
\end{solution}

\begin{solution}{exo:g9:landforms:7}
(क) पूर्व की ओर। (ख) \omterm{def:g9:landforms:beach}{पुलिन} पर उत्तर-पूर्व की ओर तिरछी तीन ऊपर की चालें, हर एक के बाद
दक्षिण की ओर सीधी नीचे वापसी की एक चाल: पूर्व की ओर बढ़ता एक टेढ़ा-मेढ़ा रास्ता। (ग) एक
\omterm{def:g9:landforms:beach}{भू-जिह्वा}, जो खाड़ी के \omterm{def:g5:water-shapes-the-land:source-mouth}{मुहाने} के पार पूर्व की ओर बढ़ती है। (घ) \omterm{def:g9:landforms:beach}{भू-जिह्वा} अपने पीछे के पानी
को लहरों से बचाती है; शांत पानी में बारीक \omterm{def:g3:stones-pebbles-sand:sand}{कीचड़} बैठ सकता
है।
\end{solution}

\begin{solution}{exo:g9:landforms:8}
(क) वह हवा से और मिट्टी से कार्बन डाइऑक्साइड लेता है। (ख) कमज़ोर अम्ल चूना पत्थर
को धीरे-धीरे घोलता है, पर ग्रेनाइट पर मुश्किल से असर करता है। (ग) पानी \omterm{def:g3:stones-pebbles-sand:rock}{चट्टान} में केवल
उसकी दरारों से और उसकी परतों के साथ-साथ घुस सकता है; वह सबसे पहले इन्हीं रास्तों को चौड़ा
करता है। (घ) सतह के नीचे की \omterm{def:g3:stones-pebbles-sand:rock}{चट्टान} घुल जाती है और ज़मीन धीरे-धीरे धँसती है; या किसी \omterm{def:g9:landforms:karst}{गुफा}
की छत ढह जाती है।
\end{solution}

\begin{solution}{exo:g9:landforms:9}
(क) जब बूँद \omterm{def:g9:landforms:karst}{गुफा} की हवा से मिलती है, तो वह अपनी कुछ कार्बन डाइऑक्साइड छोड़ देती है और
घुले हुए सारे चूना पत्थर को अब नहीं थाम पाती, जो कैल्साइट के रूप में बाहर आ जाता है।
(ख) निश्चुताश्म उन बूँदों से बढ़ता है जो फ़र्श पर छपककर फैल जाती हैं; उसकी नोक पर कोई
लटकी बूँद नहीं होती जो नली बनाए।
(ग) $\qty{30}{cm} = \qty{300}{mm}$; $300 \div 0.1 = 3000 = 3 \times 10^3$
वर्ष।
\end{solution}

\begin{solution}{exo:g9:landforms:10}
(क) $4600 \div 10^7 = 4.6 \times 10^{-4}$~m प्रति वर्ष, या \qty{0.46}{mm} प्रति
वर्ष। (ख) $0.46 - 0.2 = \qty{0.26}{mm}$ प्रति वर्ष। (ग) $\qty{1}{m} =
\qty{1000}{mm}$; $1000 \div 0.46 \approx 2200$, लगभग दो हज़ार गुना
तेज़। (घ) वह लाखों वर्षों तक बना रहता है: $\qty{0.26}{mm} \times
10^7 = 2.6 \times 10^6$~mm, लगभग \qty{2.6}{km}।
\end{solution}

\begin{solution}{exo:g9:landforms:11}
(क) आकार: अर्धचंद्राकार \omterm{def:g9:landforms:dune}{टिब्बे}, जिनका \omterm{def:g9:landforms:dune}{फिसलन फलक} दक्षिण-पूर्व की ओर है।
कारक और प्रक्रिया: उत्तर-पश्चिम से चलती हवा, \omterm{def:g5:water-shapes-the-land:deposition}{निक्षेपण}।
पदार्थ: ढीली \omterm{def:g3:stones-pebbles-sand:sand}{रेत}। (ख) आकार: भाटे पर खुलने वाले एक समतल चबूतरे के ऊपर एक \omterm{def:g9:landforms:cliff}{कगार}।
कारक और प्रक्रिया: \omterm{def:g2:water-on-earth:ocean}{समुद्र}, \omterm{def:g5:water-shapes-the-land:erosion}{अपरदन}; चबूतरा पीछे हटती \omterm{def:g9:landforms:cliff}{कगार} का छोड़ा हुआ तरंग-घर्षित
चबूतरा है। पदार्थ: सफ़ेद \omterm{def:g3:stones-pebbles-sand:rock}{चट्टान}, शायद खड़िया, एक नरम \omterm{def:g3:stones-pebbles-sand:rock}{चट्टान}। (ग) आकार: एक छेद में
ग़ायब होती धारा, दूर एक \omterm{def:g6:water-underground:spring}{सोता}। कारक और प्रक्रिया: \omterm{def:g3:stones-pebbles-sand:rock}{चट्टान} को घोलता और ज़मीन के नीचे
बहता पानी। पदार्थ: चूना पत्थर; पठार एक \omterm{def:g9:landforms:karst}{कार्स्ट} है।
\end{solution}

\begin{solution}{exo:g9:landforms:12}
(क) छोटा \omterm{def:g9:landforms:dune}{टिब्बा}: $45 \div 6 = 7.5$ वर्ष; बड़ा: $45 \div 3 =
15$ वर्ष। (ख) $45 \div 100 = 0.45$ वर्ष, लगभग पाँच या छह महीने।
(ग) उदाहरण के लिए, \omterm{def:g9:landforms:dune}{टिब्बों} पर घास लगाना: उनकी जड़ें \omterm{def:g3:stones-pebbles-sand:sand}{रेत} को थामे रखती हैं, इसलिए
हवा उसे अब नहीं ले जा पाती।
\end{solution}

\begin{solution}{pb:g9:landforms:1}
\textbf{1.} $46.0 - 30.4 = \qty{15.6}{m}$।

\textbf{2.} $15.6 \div 20 = \qty{0.78}{m}$ प्रति वर्ष।

\textbf{3.} 2012 और 2014 के बीच: दो वर्षों में $41.0 - 36.2 = \qty{4.8}{m}$,
या $4.8 \div 2 = \qty{2.4}{m}$ प्रति वर्ष।

\textbf{4.} $2.4 \div 0.78 \approx 3$ गुना तेज़। शायद एक तूफ़ानी जाड़ा, जिसमें
बड़े लटके हुए हिस्से गिरे।

\textbf{5.} \qty{78}{cm} प्रति वर्ष; $0.78 \times 100 = \qty{78}{m}$ प्रति
सदी।

\textbf{6.} हाँ: \qty{0.78}{m} प्रति वर्ष एक मीटर प्रति वर्ष से थोड़ा कम है।

\textbf{7.} $17 \div 0.78 \approx 22$ वर्ष।

\textbf{8.} \omterm{def:g9:landforms:cliff}{कगार} छलाँगों में पीछे हटती है: दो वर्ष के अंतराल में एक बड़ा गिराव हो सकता है,
या कोई भी नहीं। 20 वर्षों में छलाँगें और शांत वर्ष औसत होकर बराबर हो
जाते हैं।

\textbf{9.} किनारे को $30.4 - 10 = \qty{20.4}{m}$ और चलना है: $20.4 \div 0.78
\approx 26$ वर्ष, यानी 2050 में।

\textbf{10.} $2 \times 0.78 = \qty{1.56}{m}$ प्रति वर्ष; 2024 के $30.4 \div 1.56
\approx 19.5$ वर्ष बाद, यानी 2043 में।

\textbf{11.} \omterm{def:g9:landforms:cliff}{कगार} का किनारा किसी भी पल धँस सकता है; उस पर या उसके पास खड़ा होना
ख़तरनाक है।

\textbf{12.} 2024 के $30.4 \div 0.78 \approx 39$ वर्ष बाद:
\textbf{2063} में।
\end{solution}
